\documentclass[aspectratio=169, t, envcountsect, smaller, notheorems]{beamer}
\usetheme[sectionpage=simple, titleformat=allsmallcaps]{metropolis}
\usepackage[ngerman]{babel}
\usepackage{amssymb, mathtools, cancel, pgfplots, nicefrac}
\usepackage[ruled, linesnumbered, vlined, ngerman, algosection]{algorithm2e}
\usepackage{hyperref}

\title[Algorithmen und Datenstrukturen]{Tutorium \\ Algorithmen und Datenstrukturen}
\subtitle{Sommersemester 2019}
\author{David Luis Wiegandt}
\institute{Wissensmanagement in der Bioinformatik\\Institut für Informatik\\Humboldt-Universität zu Berlin}
\date{}
\titlegraphic{\includegraphics[width=70pt, height=70pt]{images/husiegel_bw.eps}}

\setsansfont[Path = fonts/, UprightFont = *-Regular, BoldFont = *-Medium, ItalicFont = *-Italic]{FiraSans}

\resetcounteronoverlays{algocf}
\resetcounteronoverlays{theorem}

\uselanguage{german}
\languagepath{german}
\deftranslation[to=german]{Solution}{L\"osung}
\deftranslation[to=german]{Example}{Beispiel}
\deftranslation[to=german]{Task}{Aufgabe}
\deftranslation[to=german]{Remark}{Bemerkung}

\definecolor{hublue}{RGB}{0, 55, 108}
\definecolor{lighthublue}{RGB}{9, 80, 159}
\definecolor{tealblue}{RGB}{56, 134, 151}
\definecolor{madderlake}{RGB}{204, 41, 54}
\definecolor{ashgrey}{RGB}{174, 195, 176}
\definecolor{mandarin}{RGB}{239, 123, 69}

\newcommand<>{\teal}[1]{\textit#2{\color#2{teal}#1}}
\newcommand<>{\red}[1]{\textit#2{\color#2{madderlake}#1}}
\newcommand<>{\orange}[1]{\textit#2{\color#2{mandarin}#1}}

\makeatletter
\patchcmd{\beamer@sectionintoc}{\vfill}{\vskip\itemsep}{}{}

\setbeamertemplate{frame footer}{\insertshorttitle}
\setbeamertemplate{itemize item}{\fontsize{3}{4}\raise2.5pt\hbox{\donotcoloroutermaths$\blacksquare$}}
\setbeamertemplate{itemize subitem}{\fontsize{2.5}{3.5}\raise2.5pt\hbox{\donotcoloroutermaths$\blacksquare$}}
\setbeamertemplate{itemize/enumerate body begin}{\normalsize}
\setbeamertemplate{itemize/enumerate subbody begin}{\small}
\setbeamertemplate{section in toc}[sections numbered]
\setbeamertemplate{theorems}[numbered]

\setbeamercolor{alerted text}{fg=lighthublue}
\setbeamercolor{block title example}{fg=hublue}
\setbeamercolor{footline}{fg=hublue}
\setbeamercolor{frametitle}{fg=hublue, bg=black!2}
\setbeamercolor{normal text}{fg=hublue}
\setbeamercolor{progress bar}{fg=hublue}
\setbeamercolor{title page header}{fg=white, bg=hublue}

\setbeamerfont{alerted text}{shape=\itshape}
\setbeamerfont{author}{series=\bfseries, size=\small}
\setbeamerfont{footline}{shape=\itshape}
\setbeamerfont{page number in head/foot}{size=\scriptsize}
\setbeamerfont{section title}{size=\LARGE}
\setbeamerfont{title}{size=\LARGE}

\setbeamertemplate{title page}{
\vspace*{-1pt}
\begin{beamercolorbox}[wd=\paperwidth,ht=0.618\paperheight]{title page header}
    \vfill
    \par
    \raggedright%
    \leftskip28pt%
    {\usebeamerfont{title}\inserttitle}
    \par
    \vspace*{3pt}
    {\usebeamerfont{subtitle}\insertsubtitle}
    \vspace*{-5pt}
    \vfill
    \tikz[remember picture, overlay] \node[anchor=center] at (0.8\paperwidth,0) {\inserttitlegraphic};
\end{beamercolorbox}
\vspace*{2pt}
{\usebeamerfont{author}\beamer@shortauthor}

\medskip
\fontsize{10}{8}\selectfont
\par

{\usebeamerfont{institute}\insertinstitute}
\vfill
}
\makeatother

% default styles for plots
\pgfplotsset{width=\textwidth, compat=1.9}
\pgfplotscreateplotcyclelist{hubcols}{
    {tealblue},
    {madderlake},
    {ashgrey}
}
\pgfplotsset{
    tikzDefaults/.style={
        cycle list name=hubcols,
        font=\small,
        line width = 1pt,
        legend style={
            at={(0.5,-0.2)},
            anchor=north, 
            column sep=7pt,
            draw=none,
            fill=black!2
        },
        legend columns=2,
        axis x line=bottom,
        axis y line=middle,
        tick style={draw=hublue, semithick},
        tick align=outside,
        every axis/.append style={font=\tiny},
    },
    samples=100
}

% remove rules from algos while retaining the ruled layout
\setlength{\algoheightrule}{0pt}
\setlength{\algotitleheightrule}{0pt}

% general math
\newcommand{\nat}{\mathbb{N}}
\newcommand{\natpos}{\mathbb{N}_{\geq 1}}
\newcommand{\real}{\mathbb{R}}
\newcommand{\realpos}{\mathbb{R}_{\geq 0}}
\newcommand{\powerset}{\mathcal{P}}
\newcommand{\bigO}{\mathcal{O}}
\DeclareMathOperator{\cupdot}{\mathbin{\mathaccent\cdot\cup}}
\DeclareMathOperator{\subsetfin}{\mathbin{\subset_{\mathsf{e}}}}
\newcommand{\defiff}{\mathrel{\vcentcolon\Leftrightarrow}}
\DeclarePairedDelimiter\norm{\lVert}{\rVert}

% query functions
\DeclareMathOperator{\adom}{adom}
\DeclareMathOperator{\vars}{var}
\DeclareMathOperator{\cons}{adom}
\DeclareMathOperator{\free}{frei}
\newcommand{\qfn}[1]{[\![#1]\!]}

% algorithms
\SetKwInOut{Input}{input}
\SetKwInOut{Output}{output}
\ResetInOut{output}
\SetKw{True}{true}
\SetKw{False}{false}

% theorem environments
\theoremstyle{definition}
\newtheorem{theorem}[algocf]{\translate{Theorem}}
\newtheorem{lemma}[algocf]{\translate{Lemma}}
\newtheorem{definition}[algocf]{\translate{Definition}}
\newtheorem{example}[algocf]{\translate{Example}}
\newtheorem{task}[algocf]{\translate{Task}}
\newtheorem*{solution}{\translate{Solution}}
\newtheorem*{remark}{Bemerkung}

% enumerate environments
\newenvironment{subtasks}{
    \setbeamertemplate{enumerate item}{(\alph{enumi})}
    \setbeamerfont{enumerate item}{series=\bfseries}
    \begin{enumerate}
}{
    \end{enumerate}
}

\begin{document}
    \maketitle
    
    \begin{frame}{Organisatorisches}
    \begin{block}{Zeit \& Ort}
        \begin{itemize}
            \item Mittwoch (RUD 25, 4.112) und Donnerstag (RUD 25, 4.113) von 13 bis 15 Uhr
            \item Teilnahme ist \alert{freiwillig} und \alert{flexibel}
            \begin{itemize}
                \item Ihr könnt stets an \alert{einem}, \alert{beiden} oder \alert{keinem} der Termine erscheinen
            \end{itemize}
        \end{itemize}
    \end{block}
    
    \begin{block}{Inhalte}
        \begin{itemize}
            \item Fokus ist die \alert{Wiederholung} von Inhalten aus Vorlesung und Übung
            \item Dar\"uber hinaus ist auch die \alert{Vertiefung} von Inhalten aus der Vorlesung m\"oglich
        \end{itemize}
    \end{block}
    
    \begin{block}{Bei Fragen}
        \begin{itemize}
            \item Im Tutorium oder per Mail an \alert{\href{mailto:wigandtd@hu-berlin.de}{wigandtd@hu-berlin.de}}
        \end{itemize}
    \end{block}
    \end{frame}
    
    \begin{frame}{Behandelte Themen}
        \tableofcontents
    \end{frame}
    
    \section{Laufzeitanalyse}

\begin{frame}{Hintergrund}
\begin{block}{Knappheit von Ressourcen}
    \begin{itemize}
        \item Speicherplatz (Register, RAM, HDD/SSD, Tape, ...)
        \item \alert{Rechenleistung} (Taktfrequenz)
    \end{itemize}
\end{block}
\begin{block}{Eine Problemstellung -- viele Lösungsansätze}
    \begin{itemize}
        \item Welchen Ansatz sollen wir nutzen?
        \item Ziel: Klassifiziere Lösungsansätze nach ihrer \alert{Eignung}
        \begin{itemize}
            \item Eignung $\approx$ Ressourcenbedarf
        \end{itemize}
    \end{itemize}
\end{block}
\begin{block}{Berechnungsmodell}
    \begin{itemize}
        \item Wir analysieren \alert{abstrakte Algorithmen} -- folglich st\"utzen wir unsere Analysen auf \alert{abstrakte Maschinen}
    \end{itemize}
\end{block}
\end{frame}

% \begin{frame}{Berechnungsmodell}
% Wir analysieren \alert{abstrakte Algorithmen} -- folglich st\"utzen wir unsere Analysen auf \alert{abstrakte Maschinen}.

% \medskip

% \begin{columns}[T,onlytextwidth]
% \begin{column}{0.54\textwidth}
% Beispielsweise ist die \alert{Random Access Machine} eine zur \alert{Turingmaschine} \"aquivalente und f\"ur unsere Zwecke leichter zu handhabende Abstraktion einen realen Computers.
% \end{column}
% \begin{column}{0.4\textwidth}
% \begin{tikzpicture}
% \foreach \i in {1,...,4}
% \draw[draw=hublue] (0,3.5-\i*0.5) rectangle (0.8,4-\i*0.5) node[pos=.5] {$I_{\i}$};
% \draw[draw=none, anchor=mid] (0,0.5) rectangle (0.8,1.5) node[pos=.5] {$\vdots$};
% \draw[draw=hublue] (0,0) rectangle (0.8,0.5) node[pos=.5] {$I_m$};

% \foreach \i in {1,...,6}
% \draw[draw=hublue] (1+\i*0.5,3) rectangle (1.5+\i*0.5,3.5) node[pos=.5] {$r_{\i}$};
% \draw[draw=none] (4.5,3) rectangle (5,3.5) node[pos=.5] {$\dots$};
% \end{tikzpicture}
% \end{column}
% \end{columns}
% \end{frame}

\begin{frame}<1-2>[label=first_example]{Ein einfaches Beispiel}
\begin{example}
Gegeben sei eine Liste natürlicher Zahlen $a_1, \dots, a_n \in \nat$.

Maximiere $x = a_i \cdot a_j$ mit $1 \leq i < j \leq n$.

\begin{quote}
    In Worten: Finde das \alert{gr\"o{\ss}te} Produkt, das sich aus zwei \alert{verschiedenen} Elementen der Liste bilden lässt.
\end{quote}

\begin{block}<2->{M\"ogliche Ans\"atze}
    \begin{enumerate}
        \item Betrachte alle möglichen Paare und berechne deren Produkt
        \pause
        \item<4-> Suche die beiden gr\"o{\ss}ten Elemente $a_i, a_j$ und gib $a_i \cdot a_j$ aus
    \end{enumerate}
\end{block}
\end{example}
\end{frame}

\begin{frame}
\frametitle<1>{Algorithmus}
\frametitle<2>{Algorithmus \& Analyse}
\begin{columns}[T,onlytextwidth]
    \begin{column}{0.43\textwidth}
        \vspace{-5pt}
        \begin{algorithm}[H]
        	\caption{Maximales Produkt 1}
        	\label{alg:max_prod}
        	\DontPrintSemicolon
        	\Input{$a_1, \dots, a_n \in \nat$}
            \Output{$\max \{ a_i \cdot a_j : 1 \leq i < j \leq n \}$}
        	$x \gets 0$\;
        	\For{$i \gets 1$ \KwTo $n - 1$}{
        	    \For{$j \gets i + 1$ \KwTo $n$}{
        	        \If{$a_i \cdot a_j > x$}{
        	            $x \gets a_i \cdot a_j$\;
        	        }
                }
            }
            \Return{$x$}
        \end{algorithm}
    \end{column}
    \begin{column}<2>{0.51\textwidth}
        Die \"au{\ss}ere \textbf{for}-Schleife wird $(n-1)$-mal durchlaufen, die innere \textbf{for}-Schleife $n-1, n-2, \dots, 1$-mal
        \begin{itemize}
            \item Anwendung der \alert{gaußschen Summenformel} liefert $\frac{n \, (n-1)}{2}$
        \end{itemize}
        
        Der Aufwand des \textbf{if}-Statements ist von $n$ \alert{unabh\"angig}
        \begin{itemize}
            \item Multiplikationen, Vergleiche und Zuweisungen ben\"otigen lediglich \alert{konstant} viel Zeit
        \end{itemize}
        
        Insgesamt $\leq 4 \cdot \frac{n \, (n-1)}{2} + 2$ Instruktionen
    \end{column}
\end{columns}
\end{frame}

\againframe<3->{first_example}

\begin{frame}
\frametitle<1>{Algorithmus}
\frametitle<2>{Algorithmus \& Analyse}
\begin{columns}[T,onlytextwidth]
    \begin{column}{0.44\textwidth}
        \vspace{-5pt}
        \begin{algorithm}[H]
        	\caption{Maximales Produkt 2}
        	\label{alg:max_prod_ii}
        	\DontPrintSemicolon
        	\Input{$a_1, \dots, a_n \in \nat$}
            \Output{$\max \{ a_i \cdot a_j : 1 \leq i < j \leq n \}$}
        	$k \gets 0, \quad \ell \gets 0$\;
        	\For{$i \gets 1$ \KwTo $n$}{
        	    \If{$a_i > a_k$}{
        	        $\ell \gets k, \quad k \gets i$\;
    	        }
    	        \ElseIf{$a_i > a_{\ell}$}{
    	            $\ell \gets i$\;
    	        }
            }
            \Return{$a_k \cdot a_{\ell}$}
        \end{algorithm}
    \end{column}
    \begin{column}<2>{0.5\textwidth}
        Die \textbf{for}-Schleife wird genau $n$-mal durchlaufen
        
        \medskip
        
        Innerhalb der \textbf{if}-Statements f\"uhren wir lediglich grundlegende Operationen wie Vergleiche und Zuweisungen durch
        
        \medskip
        
        Insgesamt $\leq 3 \, n + 4$ Instruktionen
    \end{column}
\end{columns}
\end{frame}

\begin{frame}<1-2>[label=second_example]{Ein komplexeres Beispiel}
\begin{example}
Wir nennen eine Sequenz $A = a_1, \dots, a_n \in \nat$ \alert{partitionierbar}, falls  es $k, \ell \in \{1, \dots, n\}$ gibt, sodass $\sum_{i=k}^{\ell} a_i = \frac{1}{2} \cdot  \sum_{i=1}^{n} a_i$.

Entscheide f\"ur eine Sequenz $A = a_1, \dots, a_n$ mit $\{ a_1, \dots, a_n \} = \{1, \dots, n\}$, ob diese partitionierbar ist.

\begin{quote}
    In Worten: Entscheide f\"ur eine \alert{Permutation} der nat\"urlichen Zahlen von $1$ bis $n$, ob es eine \alert{zusammenh\"angende Subsequenz} gibt, deren Summe die \alert{H\"alfte} der Summe der Gesamtsequenz ist.
\end{quote}

\begin{block}<2->{M\"ogliche Ans\"atze}
    \begin{enumerate}
        \item Berechne die Summe aller möglichen Subsequenzen \emph{from scratch}
        \pause
        \item<4-> Nutze \alert{dynamische Programmierung}, um fr\"uhere Teilergebnisse wiederzuverwenden
    \end{enumerate}
\end{block}
\end{example}
\end{frame}

\begin{frame}
\frametitle<1->{Algorithmus}
\frametitle<2->{Algorithmus \& Analyse}
\begin{columns}[T,onlytextwidth]
    \begin{column}{0.43\textwidth}
        \vspace{-5pt}
        \begin{algorithm}[H]
        	\caption{Subsequenz 1}
        	\label{alg:subseq_part_1}
        	\DontPrintSemicolon
        	\Input{$A = a_1, \dots, a_n \in \nat$}
            \Output{Ist $A$ partitionierbar?}
        	{\color<2>{madderlake}$s \gets \frac{n \, (n+1)}{4}$\;}
        	\For{$k \gets 1$ \KwTo $n$}{
        	    \For{$\ell \gets k$ \KwTo $n$}{
        	        {\color<3>{madderlake}$s' \gets 0$\;}
        	        \For{$i \gets k$ \KwTo $\ell$}{
        	            {\color<5>{madderlake}$s' \gets s + a_i$\;}
                    }
                    \If{{\color<3>{madderlake}$s' = s$}}{
                        {\color<4>{madderlake}\Return{\True}}
                    }
                }
            }
            {\color<4>{madderlake}\Return{\False}}
        \end{algorithm}
    \end{column}
    \begin{column}<2->{0.51\textwidth}
        \only<2->{
            Die \red<2>{Initialisierung} ben\"otigt 4 Instruktionen
            \begin{itemize}
                \item Eine Addition, eine Multiplikation, eine Division und eine Zuweisung
            \end{itemize}
        }
        
        \only<3->{
            Die \red<3>{Zuweisung} in Zeile $4$ und der \red<3>{Vergleich} in Zeile $6$ werden h\"ochstens $\frac{n \, (n+1)}{2}$-mal ausgef\"uhrt
        }
        
        \medskip
        
        \only<4->{
            Einmalig wird ein Wert \red<4>{zurückgegeben}
        }
        
        \medskip
        
        \only<5->{
            Die \red<5>{innere \textbf{for}-Schleife} wird $\frac{n \, (n+1) \, (n+2)}{6}$-mal durchlaufen (Details: n\"achste Folie!)
            \begin{itemize}
                \item Entsprechend wird auch Zeile $6$ so oft ausgef\"uhrt
            \end{itemize}
        }
        
        \only<6->{
            Insgesamt $\leq \nicefrac{1}{3} \, n \, (n+1) \, (n+5) + 5$ Instruktionen
        }
    \end{column}
\end{columns}
\end{frame}

\begin{frame}{Details zur Laufzeitschranke}
    \begin{block}{Beobachtungen}
        \begin{itemize}
            \item Die innere \textbf{for}-Schleife wird stets $(\ell - k + 1)$-mal durchlaufen
            \begin{itemize}
                \item F\"ur $k = 1$ also $1, \dots, n$-mal
                \item F\"ur $k = 2$ also $1, \dots, (n-1)$-mal
                \item F\"ur $k = n$ also $1, \dots, (n-n+1)$-mal
            \end{itemize}
            \item Als Summe ergibt sich $\sum_{i=1}^n \frac{i \, (i+1)}{2} = \frac{1}{2} \, \sum_{i=1}^n i \, (i+1) = \frac{1}{2} \, \left(\sum_{i=1}^n i^2 + \sum_{i=1}^n i \right)$
            \begin{itemize}
                \item Wir wissen, dass $\sum_{i=1}^n i = \frac{1}{2} \, n \, (n+1)$
                \item F\"ur $\sum_{i=1}^n i^2$ k\"onnen wir per vollst\"andiger Induktion zeigen, dass $\sum_{i=1}^n i^2 = \frac{1}{6} \, n \, (n + 1) \, (2 \, n + 1)$ gilt
            \end{itemize}
        \end{itemize}
    \end{block}
\end{frame}

\begin{frame}{Details zur Laufzeitschranke (2)}
    Insgesamt erhalten wir:
    
    \begin{align*}
        \frac{1}{2} \, \left(\sum_{i=1}^n i^2 + \sum_{i=1}^n i \right) &= \frac{1}{2} \, \left(\frac{1}{6} \, n \, (n + 1) \, (2 \, n + 1) + \frac{1}{2} \, n \, (n+1) \right) \\
        &= \frac{1}{2} \cdot \frac{1}{6} \, \left( n \, (n+1) \, (2\,n+1) + 3 \, n \, (n+1) \right) \\
        &= \frac{1}{2} \cdot \frac{1}{6} \, (2\,n+1+3) \, n \, (n+1) \\
        &= \frac{1}{2} \cdot \frac{1}{6} \, 2\, (n+2) \, n \, (n+1) \\
        &= \frac{1}{6} \, n \, (n+1) \, (n+2)
    \end{align*}
\end{frame}

\againframe<3->{second_example}

\begin{frame}{Dynamische Programmierung}
\begin{block}{Grundidee}
    \begin{itemize}
        \item Nutze \alert{Teilergebnisse}, um Ergebnisse f\"ur gr\"o{\ss}ere Instanzen zu bestimmen
        \item Zum Speichern von Teilergebnissen ben\"otigen wir \alert{zus\"atzlichen Speicherplatz}
    \end{itemize}
\end{block}

\begin{example}
Sei $a_1, \dots, a_n \in \nat$ eine Sequenz nat\"urlicher Zahlen.

Um f\"ur alle $1 \leq k \leq n$ die \alert{kumulierte Summe} $s_k = \sum_{i=1}^{k} a_i$ zu bestimmen, k\"onnen wir
\begin{enumerate}
    \item f\"ur jedes $k$ auf's Neue $s_k = a_1 + \cdots + a_{k}$ berechnen, oder
    \item $s_k$ aus $s_{k-1}$ bestimmen, da $s_k = s_{k-1} + a_k$.
\end{enumerate}
\end{example}
\end{frame}

\begin{frame}
\frametitle<1->{Algorithmus}
\frametitle<2->{Algorithmus \& Analyse}
\begin{columns}[T,onlytextwidth]
    \begin{column}{0.43\textwidth}
        \vspace{-5pt}
        \begin{algorithm}[H]
        	\caption{Subsequenz 2}
        	\label{alg:subseq_part_2}
        	\DontPrintSemicolon
        	\Input{$A = a_1, \dots, a_n \in \nat$}
            \Output{Ist $A$ partitionierbar?}
            {\color<2>{madderlake}
                $s \gets \frac{n \, (n+1)}{4}$\;
            	\red<5>{Sei $(m_{i,j}) \in \nat^{n \times n}$ eine Matrix\;}
        	}
        	\For{$k \gets 1$ \KwTo $n$}{
        	    \For{$\ell \gets k$ \KwTo $n$}{
        	        {\color<3>{madderlake}
            	        $m_{i,j} \gets m_{i, j - 1} + a_{\ell}$\;
            	        \If{$m_{i,j} = s$}{
            	            \red<4>{\Return{\True}}
            	        }
        	        }
                }
            }
            \red<4>{\Return{\False}}
        \end{algorithm}
    \end{column}
    \begin{column}<2->{0.51\textwidth}
        Die \red<2>{Initialisierung} ben\"otigt 4 Instruktionen
        \begin{itemize}
            \item Eine Addition, eine Multiplikation, eine Division und eine Zuweisung
        \end{itemize}
    
        \only<3->{
            Gem\"a{\ss} der gaußschen Summenformel wird der \red<3>{Rumpf} der inneren \textbf{for}-Schleife h\"ochstens $\frac{n \, (n+1)}{2}$-mal ausgef\"uhrt
            
            \begin{itemize}
                \item Dieser beinhaltet $4$ Instruktionen -- eine Subtraktion (in $m_{i, j-1}$), eine Addition, eine Zuweisung und ein Vergleich
            \end{itemize}
        }
        
        \only<4->{
            Einmalig wird ein Wert \red<4>{zur\"uckgegeben}
            
            \medskip
        }
        
        \only<5>{
            Insgesamt $2 \, n \, (n+1) + 5$ Instruktionen
            \begin{itemize}
                \item Und \red<5>{$\approx n^2$ zus\"atzlicher Speicherplatz}
            \end{itemize}
        }
        
    \end{column}
\end{columns}
\end{frame}
    %\section{\uppercase{$\bigO$}-Notation}

\begin{frame}{Definition}
\begin{columns}[T,onlytextwidth]
\begin{column}{0.47\textwidth}
    \begin{definition}[$\bigO$-Notation]
    Seien $\teal{f}, \red{g}: \nat \to \realpos$ Funktionen.
    Wir schreiben $\teal{f} \in \bigO(\red{g})$ genau dann, wenn es eine Konstante $\orange{c} \in \real_{> 0}$ gibt, sodass ab einem Index $n_0$ f\"ur alle $n \geq n_0$ gilt, dass $\teal{f}(n) \leq \orange{c} \cdot \red{g}(n)$.
    \end{definition}
    
    \medskip
    
    {\footnotesize Das hei{\ss}t, ab einem gewissen Punkt $n_0$ gibt es keinen \alert{wesentlichen} Unterschied mehr zwischen den Funktionswerten von $\teal{f}$ und $\red{g}$.}
\end{column}
\begin{column}{0.47\textwidth}
    \begin{tikzpicture}
        \begin{axis}[tikzDefaults, xmin=-4, xmax=12, ymin=0, ymax=500]
            \foreach \k in {1, ..., 8}{
                \addplot+[domain=-4:12, draw=mandarin!40, line width=0.5pt, forget plot]{1/2 * \k * x^2};
            }
            \addplot+[domain=-4:12, draw=teal]{2*x^2 - 7*x + 18};
            \addlegendentry{$2 \, x^2 - 7 \, x + 18$}
            \addplot+[domain=-4:12, draw=madderlake]{x^2};
            \addlegendentry{$x^2$}
        \end{axis}
    \end{tikzpicture}
\end{column}
\end{columns}
\end{frame}

\begin{frame}{Beispiel}
\begin{example}
Sei \teal{$f(n) \coloneqq 3 \, n + 5$} und \red{$g(n) \coloneqq 2 \, n$}.

Um zu zeigen, dass $\teal{f} \in \bigO(\red{g})$ ist, w\"ahlen wir \orange{$c=2$} und erhalten
\begin{equation*}
    \teal{\only<1>{3 \, n} \only<2->{\cancel{3 \, n}} + 5} \leq \orange{2} \cdot \red{2 \, n} = \only<1>{3 \, n} \only<2->{\cancel{3 \, n}} + n.
\end{equation*}
\only<2->{F\"ur $n \geq 5$ ist also $\teal{f(n)} \leq \orange{2} \cdot \red{g(n)}$ und somit $n_0 = 5$.}
\end{example}
\end{frame}

\begin{frame}{Bemerkungen}
\begin{remark}<+->
In der Praxis schreibt man h\"aufig $f(n) \in \bigO(g(n))$ anstelle von $f \in \bigO(g)$.

Beispielsweise ist $$\teal{2 \, n^2 + n} \in \bigO(\red{n^2})$$ \"aquivalent zu $$\teal{f} \in \bigO(\red{g}) \text{ f\"ur } \teal{f(n) \coloneqq 2 \, n^2 + n} \text{ und } \red{g(n) \coloneqq n^2}.$$
\end{remark}

\begin{remark}<+->
F\"ur Algorithmus \ref{alg:max_prod} konnten wir zeigen, dass dieser h\"ochstens $4 \cdot \frac{n \, (n-1)}{2} + 2 \in \bigO(n^2)$ Instruktionen ausf\"uhrt.
Demgegen\"uber ben\"otigt Algorithmus \ref{alg:max_prod_ii} nur h\"ochstens $3 \, n \in \bigO(n)$ Instruktionen und ist somit \alert{asymptotisch} effizienter als Algorithmus \ref{alg:max_prod}.
\end{remark}
\end{frame}

\begin{frame}{Eine leichte Aufgabe}
\begin{task}<1->
Sei $f(n) \coloneqq 6 \, n + 7$.

Zeigen Sie, dass $f(n) \in \bigO(n)$ ist, d.h. finden Sie $c \in \real_{> 0}, n_0 \in \natpos$, sodass f\"ur alle $n \geq n_0$ gilt, dass $f(n) \leq c \cdot n$.
\end{task}

\begin{solution}<2->
F\"ur $c = 7$ ist $6 \, n + 7 \leq 7 \, n = 6 \, n + n$ f\"ur alle $n \geq n_0 \coloneqq 7$.
\end{solution}
\end{frame}

\begin{frame}{Mehr Aufgaben}
\begin{task}
Beweisen oder widerlegen Sie:
\begin{subtasks}
    \item $12 \, n - 4 \in \bigO(n)$
    \item $7 \, n + 1 \in \bigO(n^2)$
    \item $\log(n^2) \in \bigO(\log n)$
\end{subtasks}
\end{task}
\end{frame}

\begin{frame}{Alternative Definition}
\begin{definition}[$\bigO$-Notation]
Seien $\teal{f}, \red{g}: \real \to \real$ Funktionen.
Wir schreiben $\teal{f} \in \bigO(\red{g})$ genau dann, wenn $\limsup_{n \to \infty} \frac{\teal{f(n)}}{\red{g(n)}} < \infty$ und $\red{g(n)} > 0$ f\"ur alle $n \in \nat$.
\end{definition}

Das folgende Theorem liefert ein bei dieser Charakterisierung h\"aufig sehr n\"utzliches Werkzeug.

\begin{theorem}[Satz von L'Hospital]
...
\end{theorem}
\end{frame}

\begin{frame}{Ein Beispiel zum Satz von L'Hospital}
\begin{example}

\end{example}
\end{frame}

\begin{frame}{Aufgaben zum Satz von L'Hospital}
\begin{task}

\end{task}
\end{frame}

\begin{frame}{Aufgabe zur Analyse von Algorithmen}
\begin{task}

\end{task}
\end{frame}

\end{document}
